\scanpage{41}
\mspage{16}

La donnée des \(\Sigma_i\) équivaut à la donnée d'un
\[
 \Sigma_*\in\mathcal M\qquad
 \left(\Sigma_*=\coprod\Sigma_i\right)
\]
et d'une partition de type \(I\) de \(\Sigma_*\) i.e. d'un morphisme
\[
 \Sigma_*\longrightarrow I_{\mathcal M}.
\]
Celle des \(\Sigma_{ij}\) équivaut à celle
d'un objet \(\Sigma_{**}\) \deleted{\unsure{et|de}} \unsure{au} dessus de \(I\times I\)
\deleted{\(\operatorname{Drap}_{\unsure{1}}(I)\)}
\deleted{\(\unsure{I\times I}\)}
(on pose \(\Sigma_{ij}=\varnothing\) si \(i\not\le j\),
\(\Sigma_{ii}=\Sigma_i\)).

La donnée des
\diagram{\begin{tikzcd}
 \Sigma_{ij}\arrow[r,hook]\arrow[d]&\Sigma_j\\
 \Sigma_i&
\end{tikzcd}}
équivaut à celle de deux morphismes ``source'' et ``but''
\diagram{\begin{tikzcd}
 \Sigma_{**}\arrow[r,"\underline b"]\arrow[d,"\underline s"']&\Sigma_*\\
 \Sigma_*&
\end{tikzcd}}
compatibles avec les \(\operatorname{pr}_1\) et \(\operatorname{pr}_2\)
de \(I\times I\) dans \(I\)
(\deleted{\uline{NB} \unsure{le cas}}
car la bigraduation \inserted{de type \(I\)}
de \(\Sigma_{**}\) est déduite de \deleted{\unsure{celle}}
\unsure{graduation} de type \(I\) de \(\Sigma_*\), et des
\(\underline b\) et \(\underline s\)) ---
\unsure{où} \(\underline b\) est \inserted{\unsure{loc.\(^{t}\)}}
un morphisme bord (i.e. est une
immersion bord) et \(\underline s\) est une fibration.

On a de plus \unsure{a}
\[
 \underline{\unsure{i|\iota}}:\ \Sigma_*\unsure{-|\longrightarrow}\Sigma_{**}
\]
``morphisme identité'', dont le composé avec \(\underline b\) et
\(\underline s\) est l'identité. La donnée des
\(\Sigma_{ijk}\unsure{-|\longrightarrow}\Sigma_{ik}\) revient à

\margin{40}
